\section{An exact certificate for the two barriers}\label{sec:certificate}
All finite decimals in this Section denote exact rational numbers.  We specify
trial sequences for Equation~\eqref{eq:energy-dual}, isolate every stationary
point of its two one-variable functions, and bound their values using rational
arithmetic.

\begin{proposition}\label{prop:real-upper}
For the determinants of Equation~\eqref{eq:determinant},
\begin{equation}\label{eq:certified-real-upper}
 \limsup_{N\to\infty}\left(\frac{\log|\Delta_N|}{n^2}
                    -\frac12\log2\right)
 \le -2.290939875<-2.2909.
\end{equation}
The separate bounds for the majorants in cases $\kappa=2$ and $\kappa=1$
are $-2.290939875$ and $-2.296789875$, respectively.
\end{proposition}

\subsection{Exact trial sequences}
For $\kappa=2$ take $d=10$ and $\lambda_2=0$.  For $\kappa=1$ take
$d=8$ and $\lambda_1=2.47405979$.  Write each sequence as
\begin{equation}\label{eq:certificate-sequences}
 u_k=l_k+\sum_{z}r_z z^k,\qquad l_k=0\quad(k>d).
\end{equation}
The following tables give all coefficients.  Every integer coefficient in
them is to be multiplied by $10^{-8}$.  The finite parts are listed in
increasing order of $k$:
\begin{center}
\begin{tabular}{ccr r r r r}
\toprule
$\kappa$&$u$&\multicolumn{5}{c}{finite coefficients}\\
\midrule
2&$p$&45559127&-50750856&-6578767&13970217&4786184\\
 &&-4292433&-576704&1311615&671564&-346453\\
2&$v$&-23910158&21152432&-2885110&-11558199&6485289\\
 &&1456821&-1912176&-2263524&2742210&-1162454\\
\bottomrule
\end{tabular}

\medskip
\begin{tabular}{ccr r r r}
\toprule
$\kappa$&$u$&\multicolumn{4}{c}{finite coefficients}\\
\midrule
1&$p$&11913521&-79993701&-21956443&37903579\\
 &&10744022&-2073193&570246&-24939103\\
1&$v$&-89913025&52874280&45168341&-30708629\\
 &&-19269841&33922112&9819141&-16992389\\
\bottomrule
\end{tabular}
\end{center}
Case $\kappa=1$ has no exponential terms.  The exponential terms for
$\kappa=2$ are as follows.  For a real base $z$ the displayed $a$ is the
coefficient of $z^k$.  For a nonreal base the two entries $a,b$ are the
coefficients of $\Re(z^k),\Im(z^k)$, in that order.  In
Equation~\eqref{eq:certificate-sequences} such a pair is represented by
$r_z=(a-ib)/2$ and $r_{\bar z}=\overline{r_z}$, with the factor $10^{-8}$
applied to $a,b$.  Each real base appears once and each nonreal base together
with its conjugate.
\begin{center}
\begin{tabular}{ccrr}
\toprule
$u$&base $z$&$a$&$b$\\
\midrule
$p$&$.85$&-9338452&---\\
&$.94$&-2141509&---\\
&$.7i$&-66277922&-31907569\\
&$.85i$&-1231651&6002645\\
&$.092+.92i$&-3225918&8928234\\
&$-.092+.92i$&-2105536&-9091287\\
\midrule
$v$&$-.8$&15199211&---\\
&$-.96$&4451662&---\\
&$.88$&2545398&---\\
&$.95$&-4932634&---\\
&$.984$&11618157&---\\
&$.78i$&27238714&-38447936\\
&$.9i$&-31341084&-30188786\\
&$.955i$&-6693542&11912254\\
&$.984i$&2055213&-21715849\\
\bottomrule
\end{tabular}
\end{center}
In particular these are real exponentially decaying sequences.  There are
ten $p$ tail terms and thirteen $v$ tail terms when conjugates are counted,
and their base moduli are at most $.94$ and $.984$, respectively.
All the $|l_k|$ and $|r_z|$ are less than $1$.

The convergent power series for the logarithm give the finite formulas
\begin{align}
 \|u\|_*^2
 &=\sum_{k=1}^d\frac{l_k^2+2l_k\sum_zr_zz^k}{k}
             -\sum_{z,z'}r_zr_{z'}\operatorname{Log}(1-zz'),\notag\\
 T(u,x)&=\sum_{k=1}^d\frac{l_kT_k(x)}{k}
             -\frac12\sum_zr_z\operatorname{Log}(1-2xz+z^2),\notag\\
 S(u,x)&=\sum_{k=1}^d\frac{l_kx^k}{k}
             -\sum_zr_z\operatorname{Log}(1-xz).
 \label{eq:certificate-log-formulas}
\end{align}
Here $\operatorname{Log}$ is the principal complex logarithm; all sums are
real by conjugation.  There is no branch ambiguity in these formulas.
For the norm and power terms the log arguments have positive real part.
For the cosine term, writing $x=\cos\theta$ gives
\[
 1-2xz+z^2=(1-ze^{i\theta})(1-ze^{-i\theta}).
\]
Both factors have positive real part, and their principal arguments sum
strictly inside $(-\pi,\pi)$.  The sum of their logarithms is therefore
exactly the principal logarithm of the product, including when the product
has negative real part.

For these choices put
\begin{align}
 X_\kappa(x)&=W_\kappa(x)+\lambda_\kappa D(x)
                         -2\kappa T(p,x)-S(v,x),\notag\\
 Y_\kappa(x)&=\beta\log x+\gamma\log(1-x)+2T(v,x).
 \label{eq:certificate-barriers}
\end{align}
The domain of $X_\kappa$ is $[-1,1)\setminus\{0\}$, and that of
$Y_\kappa$ is $(0,1)$.

\subsection{Derivative numerators and exhaustive root brackets}
Define the rational functions
\begin{align*}
 t_u(x)&=\sum_{k=1}^dl_kU_{k-1}(x)
               +\sum_z\frac{r_zz}{1-2xz+z^2},\\
 h_u(x)&=\sum_{k=1}^dl_kx^{k-1}
               +\sum_z\frac{r_zz}{1-xz}.
\end{align*}
Differentiation of Equation~\eqref{eq:certificate-log-formulas} yields
\begin{align}
 X_\kappa'(x)
 &=\frac{\alpha+2\gamma}{x}-\frac{2\eta}{1-x}
 -(\kappa+2\alpha+2\eta+2\gamma)\frac{x}{1+x^2}
 \notag\\
 &\quad-\frac{4\lambda_\kappa x}{(1+x^2)^2}
 -2\kappa t_p(x)-h_v(x),\notag\\
 Y_\kappa'(x)&=\frac\beta x-\frac\gamma{1-x}+2t_v(x).
 \label{eq:certificate-derivatives}
\end{align}
For a completely specified rational numerator, use
\begin{align}
 Q_X(x)&=x(1-x)(1+x^2)^{3-\kappa}
       \prod_{z\text{ in }p}(1-2xz+z^2)
       \prod_{z\text{ in }v}(1-xz),\notag\\
 Q_Y(x)&=x(1-x)\prod_{z\text{ in }v}(1-2xz+z^2),\notag\\
 A_X(x)&=Q_X(x)X_\kappa'(x),\notag\\
 A_Y(x)&=Q_Y(x)Y_\kappa'(x).
 \label{eq:certificate-numerators}
\end{align}
Empty products equal $1$.  These equations, the coefficient tables, and
$T_0=1,T_1=x$, $U_0=1,U_1=2x$, with recurrence
$V_{k+1}=2xV_k-V_{k-1}$, specify the four polynomials over $\mathbb Q$
without any numerical root calculation.
The denominators do not vanish on the indicated open domains:
$|1-xz|\ge1-|z|>0$, and
$|1-2xz+z^2|\ge(1-|z|)^2>0$ for real $|x|\le1$.
Conjugate factors in $Q_X,Q_Y$ have positive products, and real-base
factors are positive.  Hence $Q_X$ has the sign of $x$ on its domain and
$Q_Y>0$ on $(0,1)$.

Here is an exact root-count certificate.  If
$A(x)=\sum_{j=0}^{d_0}A_jx^j$ is the specified numerator, on a subinterval
$(b,c)$ form
\begin{equation}\label{eq:certificate-descartes-transform}
 (1+t)^{d_0}A\left(\frac{b+ct}{1+t}\right)
 =\sum_{h=0}^{d_0}a_ht^h,\qquad
 a_h=\sum_{j=0}^{d_0}\sum_{k=0}^j
 A_j\binom jk b^{j-k}c^k\binom{d_0-j}{h-k},
\end{equation}
where out-of-range binomial coefficients are zero.  Delete zero
coefficients and count consecutive sign changes.  The degrees and the
resulting counts, in the order of the consecutive intervals, are
\begin{center}
\begin{tabular}{cccll}
\toprule
$\kappa$&function&$d_0$&division points&sign variations\\
\midrule
2&$X$&36&$-1,0,1$&$9,9$\\
2&$Y$&24&$0,.25,.5,.75,1$&$5,2,2,6$\\
1&$X$&13&$-1,-.5,0,.5,1$&$1,1,2,1$\\
1&$Y$&9&$0,1$&$5$\\
\bottomrule
\end{tabular}
\end{center}
The sign-variation bound is Descartes' rule of signs
\cite[Book~III, p.~373]{Descartes1637}.
For completeness, the required bound follows by factoring
out positive real roots: multiplication of a real polynomial by $t-r$,
$r>0$, increases the number of variations by at least one.  To see this,
positive rescaling of the variable reduces to $r=1$, and zero initial
coefficients can be removed.  If $a_0,\ldots,a_d$ are the old coefficients
and $b_0,\ldots,b_{d+1}$ the new ones, then
\[
 \sum_{j=0}^h b_j=-a_h\quad(0\le h\le d),\qquad b_{d+1}=a_d.
\]
Each sign change between nonzero partial sums forces an intervening $b_j$
with the sign of the later partial sum.  Starting with $b_0=-a_0$,
these choices give an ordered subsequence with all the variations of the
$a_h$; the final coefficient $b_{d+1}$ has the opposite sign to the last
partial sum and supplies one further variation.  Factoring successively
therefore proves that the variation count bounds
the number of positive roots with multiplicity.  The substitution in
Equation~\eqref{eq:certificate-descartes-transform} sends $t>0$ bijectively
to $(b,c)$ and does not change multiplicities.

The following integers $m$ specify disjoint open brackets
$(m,m+2)/10^{10}$ for roots of the corresponding derivative:
\begin{center}
\begin{tabular}{ccrrr}
\toprule
$\kappa$&function&\multicolumn{3}{c}{$m$}\\
\midrule
2&$X$&-9601109148&-8942317572&-7608305633\\
&&-6503394794&-5185864065&-4015634158\\
&&-3108806646&-2067921826&-1589849496\\
&&1531948062&2072448179&3208186484\\
&&4381119427&5851354199&7269030693\\
&&8390277402&9332614564&9709786219\\
\midrule
2&$Y$&176402802&330649406&764952882\\
&&1227250753&2149465998&3048189112\\
&&4322699096&5564757994&6801929373\\
&&8031988371&8877037851&9577761832\\
&&9838463999&9972727815&9992037196\\
\midrule
1&$X$&-9917299785&-2259572153&2543808026\\
&&4437270259&6348298970&\\
\midrule
1&$Y$&532669786&2504239325&5701738806\\
&&7966939383&9454490138&\\
\bottomrule
\end{tabular}
\end{center}
For each integer in this table the exact sign check is
\begin{equation}\label{eq:certificate-bracket-sign}
 N_A(m)N_A(m+2)<0,\qquad
 N_A(s)=\sum_{j=0}^{d_0}A_j(10^{10})^{d_0-j}s^j.
\end{equation}
Equations~\eqref{eq:certificate-numerators},
\eqref{eq:certificate-descartes-transform}, and
\eqref{eq:certificate-bracket-sign} give rational addition and
multiplication recipes for every entry of the root certificate.
The numbers of brackets in each consecutive interval equal the displayed
variation counts.  The intermediate value theorem and the variation
bound therefore give exactly one simple root per bracket and no other
roots in those open intervals.  Possible roots at the division points
are harmless because all finite division points are evaluated separately
below.

\subsection{Rational logarithms and all candidate values}
Here is the rational procedure used for the value bounds.  For a positive
rational $s$, write $s=2^m y$ with $1\le y\le2$, and put
\begin{equation}\label{eq:certificate-H}
 H(y)=2\sum_{j=0}^{17}\frac1{2j+1}
               \left(\frac{y-1}{y+1}\right)^{2j+1}.
\end{equation}
The substitution for $\log s$ is $mH(2)+H(y)$.  If
$q=(y-1)/(y+1)$, its unscaled positive remainder is bounded by
\[
 0\le\log y-H(y)\le\frac{2q^{37}}{37(1-q^2)}.
\]
For a nonzero rational complex number $a+ib$, its real log part is half
the real log of $a^2+b^2$.  For its principal argument choose an octant:
rotate by $-k\pi/4$, $-4\le k\le4$, so the rotated real part is positive
and the imaginary-to-real ratio $t$ satisfies $|t|\le1/2$.  Select the
rotation for which $k\pi/4+\arctan t$ is the principal argument; at the
negative real axis use the value $\pi$.  The ratio $t$ is rational, since
the rotation can be performed, up to a positive scale, by repeated maps
$(a,b)\mapsto(a+b,b-a)$ or $(a-b,a+b)$.
With
\begin{equation}\label{eq:certificate-J}
 J(t)=\sum_{j=0}^{23}\frac{(-1)^jt^{2j+1}}{2j+1},
\end{equation}
substitute $k(J(1/2)+J(1/3))+J(t)$ for the argument.  The identity
\[
 \arctan(1/2)+\arctan(1/3)=\pi/4
\]
follows from the tangent addition
formula and the fact that both summands are positive and their sum is
less than $\pi/2$.  The alternating-series estimate gives
$|\arctan t-J(t)|\le |t|^{49}/49$.

All positive inputs used here, including squared moduli, lie between
$2^{-100}$ and $2^{100}$.  Indeed, the real arguments at the evaluation
points are between $.0007$ and $2$, and the squared moduli of all complex
arguments lie between $(.016)^4$ and $(1+.984)^4$, by the factor bounds
following Equation~\eqref{eq:certificate-numerators}.
Thus at most $100$ scaling steps are needed.  To make the error allowance
explicit, set
\[
 h_0=\frac{2(1/3)^{37}}{37(1-1/9)},\qquad
 j_0=\frac{(1/2)^{49}}{49}.
\]
A real log substitution costs at most $101h_0$, and an argument costs
at most $9j_0$.  In particular
$2\cdot101h_0+9j_0<7\cdot10^{-16}$, so $10^{-15}$ is an upper
bound for the absolute error of a complex log.  This also proves the
coarser allowance $10^{-11}$ per log.

All rational polynomial terms in Equation~\eqref{eq:certificate-log-formulas}
are evaluated exactly.  The total absolute log coefficient mass in
$\kappa\|p\|_*^2+\|v\|_*^2/2$ is at most
$2\cdot10^2+13^2/2=284.5$, in $X_\kappa$ it is less than
$2\cdot10+13+2=35$, and in $Y_\kappa$ it is less than $13+1=14$.
Their resulting substitution errors are respectively less than
$2.85\cdot10^{-13}$, $3.5\cdot10^{-14}$, and $1.4\cdot10^{-14}$,
all less than $10^{-8}$.
One may enclose each rational series term between consecutive multiples
of $10^{-40}$, add the displayed remainder bounds with outward signs,
and then propagate intervals linearly.  This additional rounding changes
the preceding allowances by less than $10^{-32}$.
For a complex weight $c$, use
$\Re(c\operatorname{Log}z)=\Re(c)\Re(\operatorname{Log}z)
-\Im(c)\Im(\operatorname{Log}z)$, retaining the sign of each weight in
interval multiplication.  This fully specifies a rational interval
calculation; no numerical logarithms are needed.

For transparency, Table~\ref{tab:certificate-values} gives all fifty
point-value upper bounds.  Its argument column is $10^{10}x$.  A row
marked $B$ is the left endpoint of the corresponding root bracket above;
a row marked $P$ is a finite division point.  Each displayed bound is a
rational upper bound rounded upwards to twelve decimal places after the
series remainder has been included.  The largest unrounded interval
width in this table is less than $2.16146856249281\cdot10^{-16}$.
\begin{longtable}{ccrrr}
\caption{Certified rational upper bounds at every candidate evaluation point.}\label{tab:certificate-values}\\
\toprule
$\kappa$&function&$10^{10}x$&type&upper bound\\
\midrule
\endfirsthead
\multicolumn{5}{c}{\tablename~\thetable\ (continued)}\\
\toprule
$\kappa$&function&$10^{10}x$&type&upper bound\\
\midrule
\endhead
\bottomrule
\endfoot
2&$X$&-9601109148&B&-0.984034048775\\
2&$X$&-8942317572&B&-0.984375363807\\
2&$X$&-7608305633&B&-0.984034052640\\
2&$X$&-6503394794&B&-0.984105522615\\
2&$X$&-5185864065&B&-0.984034053414\\
2&$X$&-4015634158&B&-0.984092061375\\
2&$X$&-3108806646&B&-0.984034037901\\
2&$X$&-2067921826&B&-0.984364031075\\
2&$X$&-1589849496&B&-0.984034038450\\
2&$X$&1531948062&B&-0.984033385315\\
2&$X$&2072448179&B&-0.984776982913\\
2&$X$&3208186484&B&-0.984034026898\\
2&$X$&4381119427&B&-0.984264010107\\
2&$X$&5851354199&B&-0.984034029391\\
2&$X$&7269030693&B&-0.984245044505\\
2&$X$&8390277402&B&-0.984034016294\\
2&$X$&9332614564&B&-0.984567630027\\
2&$X$&9709786219&B&-0.984033926884\\
2&$X$&-10000000000&P&-0.986727371546\\
2&$Y$&176402802&B&-1.608946411646\\
2&$Y$&330649406&B&-1.609949505120\\
2&$Y$&764952882&B&-1.608960295828\\
2&$Y$&1227250753&B&-1.609094597854\\
2&$Y$&2149465998&B&-1.608960426500\\
2&$Y$&3048189112&B&-1.608992193672\\
2&$Y$&4322699096&B&-1.608960428486\\
2&$Y$&5564757994&B&-1.608979509517\\
2&$Y$&6801929373&B&-1.608960430478\\
2&$Y$&8031988371&B&-1.608990548812\\
2&$Y$&8877037851&B&-1.608960429811\\
2&$Y$&9577761832&B&-1.609055802895\\
2&$Y$&9838463999&B&-1.608960412835\\
2&$Y$&9972727815&B&-1.609694899071\\
2&$Y$&9992037196&B&-1.608958123669\\
2&$Y$&2500000000&P&-1.608973617030\\
2&$Y$&5000000000&P&-1.608971701898\\
2&$Y$&7500000000&P&-1.608976908052\\
1&$X$&-9917299785&B&-2.778491531574\\
1&$X$&-2259572153&B&-1.324666731948\\
1&$X$&2543808026&B&-1.324655807046\\
1&$X$&4437270259&B&-1.352236629957\\
1&$X$&6348298970&B&-1.324666329425\\
1&$X$&-10000000000&P&-2.775818077526\\
1&$X$&-5000000000&P&-1.544057819905\\
1&$X$&5000000000&P&-1.347557680876\\
1&$Y$&532669786&B&-1.428286151250\\
1&$Y$&2504239325&B&-1.515602647362\\
1&$Y$&5701738806&B&-1.428335732372\\
1&$Y$&7966939383&B&-1.465686164672\\
1&$Y$&9454490138&B&-1.428335358167\\
\end{longtable}
The same calculation gives the following upper bounds for
$\kappa\|p\|_*^2+\tfrac12\|v\|_*^2$:
\[
 .778415976284\quad(\kappa=2),\qquad
 .931985203901\quad(\kappa=1).
\]
In particular the
rational substituted expressions, before their error allowances are
added, satisfy the following convenient stronger cutoffs:
\begin{center}
\begin{tabular}{cccc}
\toprule
$\kappa$&$\kappa\|p\|_*^2+\|v\|_*^2/2$&$X$ at listed points&$Y$ at listed points\\
\midrule
2&$.77843$&$-.98400$&$-1.60891$\\
1&$.93205$&$-1.32442$&$-1.42805$\\
\bottomrule
\end{tabular}
\end{center}
Each entry means a strict upper bound.  With the $10^{-8}$ allowance,
the weaker cutoffs we shall actually use are
\begin{equation}\label{eq:certificate-coarse-cutoffs}
\begin{array}{c|ccc}
 \kappa&\kappa\|p\|_*^2+\tfrac12\|v\|_*^2&X&Y\\\hline
 2&.77844&-.98399&-1.60890\\
 1&.9321&-1.3244&-1.4280.
\end{array}
\end{equation}
The logarithm procedure also gives $.693146<\log2<.693149$.

\subsection{From bracket values to global suprema}
A single derivative bound suffices on every whole closed bracket in the
root table.  All its points have distance greater than $.0007$ from both
$0$ and $1$.  In particular the bracket nearest $s=1$ has right endpoint
$9992037198/10^{10}$, whose distance from $1$ is exactly $.0007962802$.
For $|x|\le1$, the identity
$U_{k-1}(\cos\theta)=\sin(k\theta)/\sin\theta$ gives
$|U_{k-1}(x)|\le k$, including its endpoint limits.
Using the coefficient and base bounds above, the polynomial part of
$X_\kappa'$ is at most
$4\sum_{k=1}^{10}k+10=230$ in absolute value; its endpoint terms are
bounded by $(19/48+1/12)/.0007<685$; its $x/(1+x^2)$ term is less than
$3$; and its multiplier term is less than $4\lambda_1<10$.
Its two log-tail derivatives are bounded by
\[
 \frac{4\cdot10}{(.06)^2}+\frac{13}{.016}<11924.
\]
Consequently $|X_\kappa'|<12852<120000$ throughout every listed bracket.
For $Y_\kappa'$ the polynomial part is at most $110$, the endpoint
terms are bounded by $(7/48+1/12)/.0007<328$, and the log tails are
bounded by $26/(.016)^2=101562.5$.  Hence
$|Y_\kappa'|<102001<120000$ on every whole bracket as well.
These estimates hold for both cases, with empty tails in case $1$.

It follows from the mean value theorem that moving from a left bracket
endpoint to its stationary point increases either value by less than
\[
 120000\cdot\frac2{10^{10}}=.000024.
\]
Every interior stationary point is in one of these brackets or is a
division point.  All finite division points have been included in
Table~\ref{tab:certificate-values}: $-1$ for $X_2$;
$1/4,1/2,3/4$ for $Y_2$; and $-1,-1/2,1/2$ for $X_1$.
The remaining endpoints are singular and have limit $-\infty$:
$x\to0^\pm$ and $x\to1^-$ for $X_\kappa$, and $x\to0^+$ and
$x\to1^-$ for $Y_\kappa$.  The bounded tangent and multiplier terms do
not alter these limits.  Thus the root exhaustion, the finite endpoint
values, and the derivative estimate together certify the global suprema,
not merely the sampled values.

Finally $-1+\alpha+\gamma=-11/16$.  Apply
Proposition~\ref{prop:energy-reduction}, use
Equation~\eqref{eq:certificate-coarse-cutoffs}, and add $.000024$ for
each supremum.  Since $\log2>.693146$, the right side is strictly less
than, respectively,
\begin{align*}
 -\frac{11}{16}(.693146)+.77844-.98399-1.60890+.000048
    &=-2.290939875,\\
 -\frac{11}{16}(.693146)+.9321-1.3244-1.4280+.000048
    &=-2.296789875.
\end{align*}
Taking the larger of these constants proves
Proposition~\ref{prop:real-upper}.
